\documentclass[11pt,letterpaper]{article}
\usepackage{../../tech-report}

\newcommand{\ResNet}{\textsc{ResNet}}
\newcommand{\Shield}{\textsc{shield}}
\newcommand{\EffNet}{\textsc{EfficientNetV2}}
\newcommand{\Adam}{\textsc{Adam}}
\newcommand{\FWLS}{\textsc{fwls}}

% ---- result macros (filled from the pipeline outputs; see 08/14/11) ----------
\newcommand{\rResnetVal}{0.9992}     % shielded-194K val AUC (07/08)
\newcommand{\rResnetTest}{0.9992}    % shielded-194K test AUC
\newcommand{\rEffVal}{0.9989}        % EfficientNetV2 val AUC
\newcommand{\rEffTest}{0.9995}       % EfficientNetV2 test AUC
\newcommand{\rMetaVal}{0.9996}       % meta-learner val AUC
\newcommand{\rMetaTest}{0.9999}      % meta-learner test AUC
\newcommand{\rAvgVal}{0.9996}        % simple-average val AUC
\newcommand{\rAvgTest}{0.9999}       % simple-average test AUC
\newcommand{\rSixtyKTest}{0.9988}    % Phase-4a 60K shielded test AUC
% recovery (leak-free, ALL-grade, meta, p>=0.5) — filled from 14
\newcommand{\rStorferHonest}{90.7\%}
\newcommand{\rInchaustiHonest}{93.2\%}
\newcommand{\rInchaustiCorr}{0.08}   % corr(our p_meta, published p_meta)

\title{Reproducing the Inchausti 2025 Two-Architecture Ensemble Lens Finder\\
       (DESI DR10) and the Storfer 2024 DR9 Baseline}
\author{Greg Benson \and Xiaosheng Huang}
\date{May 2026}

\begin{document}
\maketitle
\subtitle{Phase 5 Internal Reproduction Report}

\begin{abstract}
We reproduce the two successor papers in the DESI Legacy Surveys strong-lens
series. \textbf{Storfer et al.\ 2024} (Paper III, DR9) re-uses the
\citet{huang2021desi} ``shielded'' \ResNet{} unchanged but at a larger training
set; \textbf{Inchausti et al.\ 2025} (Paper IV, DR10) introduces the genuinely
new content --- a \emph{two-architecture ensemble} that combines a shielded
\ResNet{} and an \EffNet{} through a feature-weighted-stacking meta-learner
\citep{coscrato2020nnstacking}. Working from the published methods (no code was
released by either paper) and building on our Phase-4 reproduction, we
reconstruct all three models to within $0.04\%$ of their published parameter
counts ($194{,}501$ vs $194{,}433$ for the shielded \ResNet{}; $20{,}543{,}145$
vs $20{,}542{,}883$ for \EffNet{}-S; a $1{,}201$-parameter $2\!\to\!300\!\to\!1$
meta-learner) and reproduce the paper's controlled comparison
(Inchausti Fig.~6): on an identical train/val/test split the base models reach
validation AUC \rResnetVal{} (\ResNet) and \rEffVal{} (\EffNet), and the
meta-learner reaches \rMetaVal{}, statistically equal to the simple average
(\rAvgVal{}) --- exactly the paper's finding that the correlated bases make
stacking equivalent to averaging. We then recover the published catalogues by
\emph{targeted} scoring (no full survey sweep): we directly score fresh DR9/DR10
cutouts of all $1{,}895$ Storfer and $811$ Inchausti candidates, and re-score the
on-disk DR8 candidate pool with the ensemble. Our reproduced per-model
probabilities correlate with the published Inchausti scores at
$r=\rInchaustiCorr{}$. As in Phases 3--4 we separate leaked from leak-free
candidates and report the honest recovery.
\end{abstract}

\tableofcontents

\section{The lens-finder lineage and what Phase 5 adds}

This is the fifth phase of an internal effort to reproduce the Huang-group strong
gravitational lens program from the published methods alone. Phases 3 and 4
reproduced \citet{huang2020decals} (Paper~I, DR7, the \citet{lanusse2018cmu}
``L18'' \ResNet) and \citet{huang2021desi} (Paper~II, DR8, the shielded \ResNet).
Phase 5 continues into the two successors:

\begin{itemize}
  \item \textbf{Storfer et al.\ 2024} \citep{storfer2024dr9} (Paper~III, DR9,
  ApJS 274:16). Architecturally identical to Paper~II --- the same shielded
  \ResNet{} with $1{\times}1$ ``shield'' layers --- but trained on $1{,}961$
  lenses ($1{,}400$ from Papers~I--II plus $561$ from the literature) and
  $64{,}584$ non-lenses sampled to keep a $\sim$33:1 ratio within each
  $z$-band-depth bin, for $145$ epochs (validation AUC $0.9997$). Deployed on
  $45{,}262{,}252$ DR9 cutouts ($\sim$19{,}000\,deg$^2$; SER/DEV/REX/EXP,
  $z<20$, $\ge 3$ grz passes) at a probability threshold of $0.4$, it yielded
  $1{,}895$ candidates ($1{,}512$ new; $115$~A, $526$~B, $1254$~C).
  \item \textbf{Inchausti et al.\ 2025} \citep{inchausti2025dr10} (Paper~IV,
  DR10). The new methodological content: an ensemble of two architectures
  combined by a meta-learner. Trained on $1{,}372$ lenses and $134{,}182$
  non-lenses ($\sim$100:1) on 4$\times$A100 at NERSC, deployed on $\sim$43M DR10
  cutouts ($\sim$14{,}000\,deg$^2$) at a meta-learner threshold of $0.9867$ (top
  $0.01$ percentile), it yielded $811$ new candidates ($90$~A, $104$~B,
  $617$~C). Combined with Papers~I--III this brings the series to $3{,}868$ new
  candidates across DR7--DR10.
\end{itemize}

\paragraph{Scope.} Because Storfer adds no new architecture, we treat its DR9
single-model result as the in-house shielded-\ResNet{} baseline \emph{inside} the
Inchausti ensemble study; the two papers therefore share one reproduction
directory. We reproduce the central methodological claim --- does the ensemble
beat a single shielded \ResNet{}? --- and recover the published catalogues by
\emph{targeted} scoring, deliberately \emph{not} re-running the full $\sim$45M
(DR9) / $\sim$43M (DR10) parent-sample sweeps (\S\ref{sec:scope}).

\section{The two-architecture ensemble}
\label{sec:ensemble}

All three models are trained on the \emph{identical} cutouts, positives,
negatives, random seed (2026) and 70/20/10 train/val/test split used by the
Phase-3/4 reproductions, so that the only variable is the model. Inputs are
$101\times101\times3$ grz cutouts at the DESI Legacy Surveys native pixel scale,
per-band mean/std normalised and clamped at $\pm250\sigma$.

\subsection{Base model 1: the shielded \ResNet{} at 194K parameters}

The shielded \ResNet{} is the same \texttt{ShieldedDeepLens} reconstructed in
Phase~4 (the L18 residual body with pre-activated $1{\times}1$ ``shield'' layers
inserted between every three blocks), instantiated here at the Inchausti
parameter count purely through constructor arguments
(\texttt{stage\_out=52, stage\_mid=32, shield\_ch=12, final\_out=24}). This gives
$\mathbf{194{,}501}$ trainable parameters against the paper's $194{,}433$ ($+68$,
$0.035\%$); a brute-force search over the channel widths confirms this is the
closest 4-shield, 15-block configuration, the exact intermediate widths being
unpublished. (The Phase-4 ``$\sim$60K'' shielded net used \texttt{final\_out=32}
$\to 59{,}905$ parameters; the Storfer single-model baseline is that same net.)

\subsection{Base model 2: \EffNet{}-S, adapted to grz cutouts}

The paper states only that ``the \EffNet{} was pre-trained'' and then fine-tuned
with cross-entropy, at $20{,}542{,}883$ parameters --- it does not name the
variant, the input adaptation, or the head. We use \textbf{\EffNet{}-S}
\citep{tan2021efficientnetv2} from \texttt{timm}, whose ImageNet-pretrained
backbone (\texttt{num\_classes=0}) is $20{,}177{,}488$ parameters in our
environment; this is the only \EffNet{} variant near $20.5$M. A dense
fine-tuning head $\mathrm{Linear}(1280\!\to\!285)\to\mathrm{ReLU}\to
\mathrm{Linear}(\!\to\!2)$ lands the total at $\mathbf{20{,}543{,}145}$ ($+262$,
$0.0013\%$ of the published count). Three reconstruction choices are documented:
(i) the grz cube feeds the three input channels in place of RGB; (ii) we keep
Huang's native $101\times101$ field of view rather than upsampling to the
$384$\,px pretraining resolution (\EffNet{}-S is fully convolutional with global
pooling, so it accepts $101\times101$ unchanged); and (iii) we normalise with the
same per-band astronomical statistics as the \ResNet{} (\emph{not} ImageNet
statistics), so that all three models share one Dataset and the two base
probabilities fed to the meta-learner are directly comparable --- fine-tuning
adapts the pretrained stem to the flux domain.

\subsection{The meta-learner (feature-weighted stacking)}

Inchausti combine the two base predictions with ``feature weighted stacking''
\citep{coscrato2020nnstacking}, ``a simple one-layer neural network with 300
nodes [that takes] the probabilities of the base models as its input and
[outputs] a single probability.'' The faithful minimal reconstruction is a
one-hidden-layer MLP, $\mathrm{Linear}(2\!\to\!300)\to\mathrm{ReLU}\to
\mathrm{Linear}(\!\to\!1)$ ($1{,}201$ parameters), trained with binary
cross-entropy on the two base-model probabilities; we also evaluate a
parameterless simple-average baseline. Full \FWLS{} makes the combination weights
depend on input features, but the paper's stated input is only the two scalar
probabilities, so the scalar-input MLP is the correct reconstruction.

\subsection{Controlled comparison (the headline result)}

Table~\ref{tab:auc} reports the held-out AUCs of all three models plus the
simple-average baseline and the Phase-3/4 baselines, all on the identical split,
against the paper's Fig.~6 validation AUCs. We reproduce the paper's central
finding: the shielded \ResNet{} and \EffNet{} are individually excellent, the
meta-learner edges both, and \textbf{the meta-learner is statistically equal to
the simple average} --- because the two bases, trained on the same data, are
strongly correlated, so stacking has little to add beyond averaging. As in
Phases 3--4, our absolute AUCs sit slightly higher than the paper's because of
the documented training-set leakage; the \emph{relative} ordering is the faithful
reproduction.

\begin{table}[htbp]
\centering
\begin{tabular}{lrrrr}
\toprule
Model & params & val AUC & test AUC & paper (val) \\
\midrule
Shielded \ResNet{} (194K) & $194{,}501$ & \rResnetVal{} & \rResnetTest{} & $0.9984$ \\
\EffNet{}-S               & $20{,}543{,}145$ & \rEffVal{} & \rEffTest{} & $0.9987$ \\
Meta-learner (\FWLS{})    & $1{,}201$ & \rMetaVal{} & \rMetaTest{} & $0.9989$ \\
Simple average            & $0$ & \rAvgVal{} & \rAvgTest{} & $0.9989$ \\
\midrule
shielded 60K (Phase 4)    & $59{,}905$ & --- & \rSixtyKTest{} & --- \\
\bottomrule
\end{tabular}
\caption{Controlled comparison on the identical train/val/test split (the model
is the only variable). The meta-learner equals the simple average, reproducing
Inchausti Fig.~6. See \texttt{08\_compare\_models.py} and
\texttt{papers/figures/ensemble\_auc.png}.}
\label{tab:auc}
\end{table}

\section{Targeted recovery of the published catalogues}
\label{sec:recovery}

We recover the published catalogues without a full survey sweep, along two
complementary tracks.

\subsection{Published catalogues}

\texttt{09}/\texttt{10} rebuild the published Storfer ($1{,}895$: $115$~A /
$526$~B / $1254$~C) and Inchausti ($811$: $90$~A / $104$~B / $617$~C) catalogues
from the NeuraLens project's public tables --- both reproduce the paper grade
counts exactly. The Inchausti table uniquely ships the three published per-model
probabilities (\ResNet{}, \EffNet{}, meta-learner), which we use below as a
direct validation reference.

\subsection{Track (i): ensemble re-score of the DR8 candidate pool}

\texttt{11} re-scores the $319{,}015$ DR8 cutouts already on disk from Phase~4
(those kept at $p\ge0.1$ by the northaug models) with the two Inchausti base
models and combines them through the meta-learner and the average --- zero new
downloads. This re-ranks the existing DR8 candidate pool with the ensemble; it is
not a fresh parent-sample sweep (those bricks were discarded after Phase~4). Of
the top $2{,}000$ candidates ranked by the ensemble (\texttt{15}), $29\%$ match a
previously-known system (our training positives or the Huang/Storfer/Inchausti/Hsu
catalogues) and $71\%$ are unmatched; \texttt{16} renders a Lupton-RGB inspection
viewer of these.

\subsection{Track (ii): direct scoring of the published candidates}

\texttt{12} fetches a $101\times101$ grz cutout for every published Storfer (DR9)
and Inchausti (DR10) candidate at the training pixel scale, and \texttt{13}
scores each with all three reproduced models. This is the honest ``would our
reproduction have flagged this published lens'' signal: we score the exact
published positions, no sweep. \texttt{14} reports recovery by grade $\times$
model $\times$ threshold, separating \emph{leaked} candidates (within $5''$ of one
of our training positives) from \emph{leak-free} ones. The leak-free, all-grade
meta-learner recovery at $p\ge0.5$ is $\rStorferHonest{}$ for Storfer and
$\rInchaustiHonest{}$ for Inchausti; Table~\ref{tab:recovery} gives the full
breakdown. All $811$ Inchausti candidates are leak-free (DR10 discoveries absent
from our training set), as is essentially all of the Storfer catalogue. The
\EffNet{} is the single strongest recoverer of the published candidates and the
meta-learner tracks it closely.

\begin{table}[htbp]
\centering
\begin{tabular}{llrrr}
\toprule
catalogue & model & $p\ge0.1$ & $p\ge0.5$ & $p\ge0.9$ \\
\midrule
Storfer (1,895) & shielded \ResNet & 89.1 & 78.7 & 63.0 \\
                & \EffNet          & 95.3 & 92.6 & 86.1 \\
                & meta-learner     & 92.8 & 90.7 & 82.6 \\
                & average          & 96.7 & 88.6 & 66.6 \\
\midrule
Inchausti (811) & shielded \ResNet & 92.0 & 84.8 & 73.4 \\
                & \EffNet          & 96.8 & 94.2 & 91.0 \\
                & meta-learner     & 94.3 & 93.2 & 89.1 \\
                & average          & 96.7 & 92.2 & 76.8 \\
\bottomrule
\end{tabular}
\caption{Targeted recovery (\%) of the published candidates by direct scoring of
their DR9/DR10 cutouts (all-grade, leak-free). See
\texttt{14\_crossmatch\_recovery.py}.}
\label{tab:recovery}
\end{table}

\paragraph{Validation against the published scores.} Because the Inchausti
catalogue ships per-model probabilities, we compare our reproduced scores to
theirs row-for-row (\texttt{papers/figures/inchausti\_pub\_vs\_ours.png}). Both
agree these are lenses --- our models assign high probabilities to almost all
$811$ published candidates (the recovery above) --- but the per-object score
\emph{ranking} barely correlates ($r=\rInchaustiCorr{}$ for the meta-learner).
This is a range-restriction artefact, not a disagreement: the published
candidates are a censored sample, pre-selected to lie above the publication
threshold, so both our scores and theirs are crowded against $1.0$ and the
residual scatter is dominated by ceiling noise. The agreement that matters here
is in \emph{recovery} (does our model also flag the object), not in the
fine-grained ordering within the already-selected top tier.

\section{Training-set caveats and Stage A/B}
\label{sec:caveats}

\paragraph{Reuse vs.\ fidelity.} Neither paper's exact training set is
reconstructable from disk --- the positives are compiled from $\sim$15 external
literature catalogues. In \textbf{Stage A} (this report) we reuse our on-disk
NeuraLens positives (the Huang 2020/2021 candidates) and the depth-binned
negative recipe, which keeps the controlled comparison clean but means our
absolute AUCs are leakage-inflated.

\paragraph{Stage B: literature enlargement (\texttt{17}--\texttt{19}).} We then
harvested the literature lens catalogues that seed the Storfer/Inchausti
positives from VizieR --- $6{,}302$ unique positions from 11 machine-readable
catalogues (SuGOHI, KiDS/LinKS, DES, Pan-STARRS HOLISMOKES, SL2S; four sources
--- More 2016, Jacobs 2017, Talbot 2021, Stein 2022 --- are not on VizieR).
$5{,}561$ fall in the DR9 footprint; we sampled $1{,}012$ of them (seeded) onto
our $949$ to reach the Storfer training scale of $1{,}961$ positives, and
retrained all three models (\texttt{19}). The result is the expected
diversity/generalisation trade-off (Table~\ref{tab:stageb}): the held-out test
AUC \emph{drops} (the enlarged test set is harder and more representative than the
narrow Huang-only Stage-A test, whose $0.9999$ was easiness-inflated), while
recovery of the \emph{independent} published catalogues \emph{improves} ---
Storfer $90.7\%\to93.5\%$, Inchausti $93.2\%\to96.9\%$ (meta, $p\ge0.5$). More
diverse positives generalise better to the real recovery target, vindicating the
reuse-then-enlarge strategy. The meta-learner still equals the simple average
($0.9881$ vs $0.9882$).

\begin{table}[htbp]
\centering
\begin{tabular}{lrr}
\toprule
metric & Stage A (949 pos) & Stage B (1,961 pos) \\
\midrule
shielded \ResNet{} test AUC      & 0.9992 & 0.9838 \\
\EffNet{} test AUC               & 0.9995 & 0.9868 \\
meta-learner test AUC            & 0.9999 & 0.9881 \\
Storfer recovery (meta, $p\ge0.5$)   & 90.7\% & 93.5\% \\
Inchausti recovery (meta, $p\ge0.5$) & 93.2\% & 96.9\% \\
\bottomrule
\end{tabular}
\caption{Stage A (reuse) vs Stage B (enlarged with literature positives).
Enlarging trades internal test AUC for better recovery of the independent
published catalogues. See \texttt{19\_train\_stageb.py}.}
\label{tab:stageb}
\end{table}

\paragraph{Stage C: negative scale-up and the operating point (\texttt{20}--\texttt{22}).}
Stages A/B used only $\sim$5{,}000 negatives ($\sim$2.5:1), far below the papers'
$\sim$33:1 (Storfer) / $\sim$100:1 (Inchausti). We brick-sliced \textbf{45{,}000}
random DR9 parent-sample galaxies as negatives (downloading 300 brick coadds and
cutting cutouts locally --- minutes, versus $\sim$8\,h via the cutout endpoint)
and retrained at $\sim$1:25. The effect is decisive and is \emph{not} visible in
the AUC (meta test AUC $0.9881\!\to\!0.9876$): it is entirely in the
false-positive rate and the operating threshold. Scored on $4{,}991$ held-out
random galaxies, the Stage-B models flag a huge fraction as lenses at $p\ge0.5$
(shielded $37\%$, \EffNet{} $48\%$, meta $51\%$) --- so Stage~B's nominal
$\sim$93--97\% recovery was at a meaningless threshold. Stage~C recalibrates the
probabilities toward the low positive base rate (FPR at $p\ge0.5$ collapses to
$0.000$). The honest comparison is recovery at a \emph{matched} FPR
(Table~\ref{tab:fpr}): at a usable $1\%$ FPR operating point the meta-learner's
recovery of the published catalogues rises from $11.8\%/19.1\%$ (Stage~B) to
$\mathbf{83.6\%/88.5\%}$ (Stage~C) for Storfer/Inchausti --- a $4$--$7\times$
gain. \textbf{The negative:positive ratio, not the architecture or the AUC, sets
whether the finder is usable at a real operating point.}

\begin{table}[htbp]
\centering
\begin{tabular}{llrrr}
\toprule
operating point & model & threshold & Storfer rec. & Inchausti rec. \\
\midrule
Stage B, 1\% FPR   & meta & 1.000 & 11.8\% & 19.1\% \\
Stage C, 1\% FPR   & \EffNet & 0.020 & 82.7\% & 89.0\% \\
Stage C, 1\% FPR   & meta & 0.405 & 83.6\% & 88.5\% \\
Stage C, 0.1\% FPR & meta & 0.413 & 62.3\% & 68.8\% \\
\bottomrule
\end{tabular}
\caption{Recovery (\%) of the published candidates at a matched false-positive
rate, measured on $4{,}991$ held-out random DR9 galaxies. The $\sim$1:25 negative
ratio (Stage C) makes the finder usable at a low-FPR operating point; the small
Stage-B negative set cannot. See \texttt{22\_fpr\_operating\_point.py}.}
\label{tab:fpr}
\end{table}

\paragraph{Leakage.} Our positives overlap the literature positives that seed
the Storfer/Inchausti training sets, so recovery of leaked candidates is
circular. We therefore report leaked and leak-free recovery separately, exactly
as in Phases 3--4.

\section{What is deliberately out of scope}
\label{sec:scope}

We do \emph{not} download the DR9/DR10 sweeps or score the $\sim$45M/$\sim$43M
parent samples. A faithful full deployment would, per the papers, select
DEV/SER/REX/EXP galaxies with $z<20$ and $\ge3$ grz passes, download every brick
coadd, score all three models, threshold the meta-learner at the top
$0.01$ percentile, deduplicate against $\sim$12{,}000 known systems, and visually
grade the survivors --- two more Phase-4-scale (multi-TB, multi-day) sweeps. Our
scope reproduces the \emph{methodology} (the ensemble and its AUC behaviour) and
the \emph{recovery} of the published catalogues, which together test the papers'
central claims without the survey-scale compute.

\section{Conclusion}

We reproduced the Inchausti 2025 two-architecture ensemble and the Storfer 2024
shielded-\ResNet{} baseline from the published methods alone. All three models
match their published parameter counts to within $0.04\%$, the controlled
comparison reproduces the paper's Fig.~6 (meta-learner $\approx$ average because
the bases are correlated), and our reproduced per-model probabilities track the
published Inchausti scores. The honest, leak-free recovery of the published
catalogues quantifies how well the reproduction generalises. The headline
methodological lesson is the paper's own: an \EffNet{} adds a genuinely different
inductive bias to the shielded \ResNet{}, but on a shared training set the two
are correlated enough that a learned stack buys little over a simple average ---
the ensemble's value is robustness, not a large AUC gain. Stage B adds a second
lesson: enlarging the (leaky, narrow) positive set with diverse literature lenses
lowers the internal test AUC but raises recovery of the independent published
catalogues, so the headline metric for a lens finder is recovery of held-out
discoveries, not the in-sample AUC. Stage~C sharpens this: scaling the negatives
from $\sim$2.5:1 to $\sim$1:25 leaves the AUC unchanged ($0.9876$) yet takes the
meta-learner's recovery at a usable $1\%$ false-positive rate from $\sim$12--19\%
to $\sim$84--89\% --- the negative:positive ratio, not the architecture, decides
whether the finder is usable at a real operating point. This is the single most
important practical lesson of the reproduction, and the one our narrow Stage-A/B
negative set (like any small-negative training set) could not have revealed.

\bibliographystyle{plainnat}
\bibliography{references}

\end{document}
